\bookchapter{บทที่ 34:คู่มือเชื่อมต่อรถยนต์ขั้นลึก—การผสานระหว่างโทรศัพท์และรถยนต์}

\chapterintro{รถยนต์กำลังวิวัฒนาการจาก “ล้อสี่ล้อบวกกับโซฟา” เป็น “ล้อสี่ล้อบวกกับสมาร์ทโฟน” ไม่ว่าคุณจะขับรถน้ำมันหรือรถพลังงานใหม่ ความสัมพันธ์ระหว่างสมาร์ทโฟนกับระบบรถยนต์ต่างกำหนดประสบการณ์การขับขี่ของคุณ บทนี้จะพาคุณทำความเข้าใจเส้นทางหลัก 3 ทางของการเชื่อมต่อระบบรถยนต์ คุณสมบัติของระบบรถยนต์แต่ละค่าย ตลอดจนเทคนิคเล็ก ๆ ที่ทำให้การขับขี่ราบรื่นยิ่งขึ้น}

\booksection{สามเส้นทางหลักของการเชื่อมต่อระบบรถยนต์}

ปัจจุบันการ “เชื่อมต่อ” ระหว่างสมาร์ทโฟนกับรถยนต์มีอยู่ 3 แนวคิดหลัก เข้าใจแนวคิดเหล่านี้ก็เท่ากับเข้าใจทั้งระบบนิเวศ:

\begin{itemize}
  \item \textbf{Cast หน้าจอสมาร์ทโฟนไปยังระบบรถยนต์}:ระบบรถยนต์ทำหน้าที่เป็นแค่ "จอแสดงผล" และ "ลำโพง" ส่วนการประมวลผลและข้อมูลแผนที่มาจากสมาร์ทโฟนทั้งหมด โซลูชันที่เป็นตัวแทน: Apple CarPlay, Google Android Auto, Huawei HiCar, Baidu CarLife ข้อดีคือนำระบบนิเวศของสมาร์ทโฟนมาใช้ในรถได้ทันทีโดยไม่ต้องติดตั้ง App ใหม่ ข้อเสียคือความสามารถของฮาร์ดแวร์ระบบรถยนต์เองถูกละเลย
  \item \textbf{การผสานรวมอย่างลึกซึ้งระหว่างสมาร์ทโฟนและระบบรถยนต์}:ระบบรถยนต์เป็นระบบอิสระ แต่แชร์บัญชี แอปพลิเคชัน และข้อมูลกับสมาร์ทโฟน แม้กระทั่งสามารถ "สานต่อ" งานได้ โซลูชันที่เป็นตัวแทน: ห้องโดยสาร HarmonyOS (AITO/Luxeed), BYD DiLink, NIO OS, ระบบรถยนต์ Li Auto ข้อดีคือประสบการณ์การใช้งานเป็นหนึ่งเดียว ข้อเสียคือมักถูกผูกกับระบบนิเวศของแบรนด์เดียวกัน
  \item \textbf{CarPlay 2.0 เวอร์ชันผสานรวมอย่างลึกซึ้ง}:CarPlay รุ่นใหม่ที่ Apple เปิดตัวในปี 2025 ไม่ใช่แค่ "การ Cast หน้าจอ" อีกต่อไป แต่เข้าควบคุมฟังก์ชันหลักของระบบรถยนต์ เช่น แผงหน้าปัด การควบคุมแอร์ การปรับเบาะนั่ง แบรนด์อย่าง Ferrari, Aston Martin, Porsche เป็นกลุ่มแรกที่นำไปใช้ นี่ถือเป็นสัญลักษณ์ว่าระบบรถยนต์กำลังเข้าสู่ยุคใหม่ที่ "ผู้ผลิตสมาร์ทโฟนเป็นผู้นำ"
\end{itemize}

\booksection{เปรียบเทียบ CarPlay / Android Auto / HiCar}

โซลูชันการ Cast หน้าจอหลัก 3 ทางมีจุดเน้นที่แตกต่างกัน เลือกรถหรือเลือกสมาร์ทโฟนสามารถเทียบกับตารางด้านล่างได้:

\begin{center}
\begin{tabularx}{\linewidth}{|X|X|X|X|X|}
\hline
\textbf{โซลูชัน} & \textbf{ระบบปฏิบัติการสมาร์ทโฟนที่รองรับ} & \textbf{รุ่นรถที่รองรับ} & \textbf{ความครบครันของฟังก์ชัน} & \textbf{การใช้งานในประเทศไทย} \tabularnewline
\hline
Apple CarPlay & เฉพาะ iPhone & กว้างขวางมาก รองรับรถยนต์เกือบทุกรุ่น & สูง แผนที่/เพลง/โทรศัพท์/ข้อความ ใช้ได้หมด & ดี ใช้งานได้ตามค่าเริ่มต้น \tabularnewline
Android Auto & Android (สมาร์ทโฟนในประเทศต้องติดตั้ง Service Framework) & ค่อนข้างกว้างขวาง รถยนต์ร่วมทุน/นำเข้าเป็นส่วนใหญ่ & สูง เทียบเท่า CarPlay & ปานกลาง ต้องใช้ VPN เพื่อเปิดใช้งาน \tabularnewline
Huawei HiCar & สมาร์ทโฟน Huawei เป็นหลัก รองรับ Android อื่น ๆ บางส่วน & รถยนต์ผลิตในประเทศและรถยนต์ร่วมทุนบางรุ่น & ค่อนข้างสูง เชื่อมต่อกับห้องโดยสาร HarmonyOS ได้ & ดีมาก ระบบนิเวศในประเทศ \tabularnewline
Baidu CarLife & ทั้ง iOS และ Android & ค่อนข้างกว้างขวาง พบบ่อยในการติดตั้งเพิ่มบนรถรุ่นเก่า & ปานกลาง อัปเดตช้ากว่าปกติ & ดีมาก ระบบนิเวศในประเทศ \tabularnewline
\hline
\end{tabularx}
\end{center}

\booksection{ห้องโดยสาร HarmonyOS: การส่งต่อแบบไร้รอยต่อระหว่างสมาร์ทโฟนและหน้าจอรถยนต์}

หากพูดถึง CarPlay ว่าเป็น “การส่งภาพจากโทรศัพท์ไปยังหน้าจอรถยนต์” ห้องโดยสาร HarmonyOS ก็เปรียบเสมือน “โทรศัพท์กับหน้าจอรถยนต์ที่เป็นหนึ่งเดียวกัน” การทำงานร่วมกันระหว่างโทรศัพท์ Huawei กับหน้าจอรถยนต์ Wenjie/Zhijie นั้นเป็นตัวแทนของมาตรฐานสูงสุดของการเชื่อมต่อรถยนต์ในประเทศจีนปัจจุบัน:

\begin{itemize}
  \item \textbf{ซูเปอร์เดสก์ท็อป}: หน้าจอรถยนต์สามารถเรียกใช้ App บนโทรศัพท์ได้โดยตรง โดยไม่ต้องติดตั้งใหม่บนรถยนต์ เพลงที่กำลังฟังอยู่บนโทรศัพท์ หรือเส้นทางนำทางที่ตั้งไว้ จะถูกส่งต่อไปยังหน้าจอรถยนต์โดยอัตโนมัติเมื่อขึ้นรถ และกลับมาที่โทรศัพท์เมื่อลงจากรถ ทั้งหมดเป็นกระบวนการที่ไร้รอยต่อ
  \item \textbf{การส่งต่อแอปพลิเคชัน}: เส้นทางนำทางที่ตั้งไว้บนโทรศัพท์ จะถูกหน้าจอรถยนต์รับช่วงต่อโดยอัตโนมัติเมื่อปิดประตูรถ เมื่อถึงลานจอดรถและลงจากรถ โทรศัพท์จะทำการนำทางเดินเท้าต่อไปยังจุดหมายปลายทาง "กิโลเมตรสุดท้าย" จะไม่ถูกตัดขาดอีกต่อไป
  \item \textbf{การใช้บัญชีร่วมกันระหว่างโทรศัพท์และรถยนต์}: เมื่อล็อกอินบัญชี Huawei หน้าจอรถยนต์จะซิงค์บัญชีสมาชิกเพลง รายการบุ๊คมาร์กแผนที่ ปฏิทิน และรายชื่อผู้ติดต่อของคุณโดยอัตโนมัติ หน้าจอรถยนต์และโทรศัพท์คืออุปกรณ์ที่ใช้ตัวตนเดียวกันแต่ต่างรูปแบบกัน ไม่ใช่สองระบบแยกจากกัน
\end{itemize}

\booksection{ระบบหน้าจอรถยนต์จากค่ายผู้ผลิตรถยนต์พลังงานใหม่}

หน้าจอรถยนต์ไฟฟ้าในประเทศได้พัฒนาจาก “นำทาง + เพลง” ไปเป็นอุปกรณ์อัจฉริยะเคลื่อนที่แท้จริง โดยแต่ละแบรนด์มีจุดเด่นที่แตกต่างกัน:

\begin{itemize}
  \item \textbf{NIO OS ของ NIO}: ผู้ช่วยเสียง NOMI เป็นมาตรฐานในวงการ สามารถรับคำสั่งได้หลายคำสั่ง สนทนาต่อเนื่องตามบริบทได้ และยังมีสีหน้าแสดงอารมณ์ NIO Phone ทำให้โทรศัพท์กลายเป็น "กุญแจรถ + หน้าจอที่สอง" ผูกกับระบบนิเวศของ NIO อย่างลึกซึ้ง
  \item \textbf{หน้าจอรถยนต์ Li Auto}: จุดเด่นคือการโต้ตอบแบบหลายหน้าจอ หน้าจอของผู้โดยสารข้างคนขับสามารถดูซีรีส์ได้อย่างอิสระ หน้าจอสำหรับเด็กสามารถดูการ์ตูนได้ "Task Master" ช่วยให้คุณปรับแต่งสถานการณ์ได้ เช่น "เปิดประตูแล้วเพลงจะหยุดชั่วคราว ปิดประตูแล้วเล่นต่อ"
  \item \textbf{หน้าจอรถยนต์ Xiaomi HyperOS}: เชื่อมต่อลึกซึ้งกับโทรศัพท์ Xiaomi นาฬิกา และเครื่องใช้ไฟฟ้า นาฬิกาปลุกที่ตั้งไว้บนโทรศัพท์จะแจ้งเตือนบนหน้าจอรถยนต์ เมื่อนำทางกลับบ้านด้วยรถ แอร์ที่บ้านจะเปิดทำงานล่วงหน้า HyperOS เชื่อมต่อ "ระบบนิเวศบุคคล รถยนต์ และบ้าน" เข้าด้วยกัน
  \item \textbf{DiLink ของ BYD}: จุดเด่นคือหน้าจอที่หมุนได้ ใช้แนวตั้งสำหรับดูวิดีโอสั้น แนวนอนสำหรับดูแผนที่ ปรับทิศทาง App ได้อัตโนมัติ มี App Store ในตัว ความอุดมสมบูรณ์ของระบบนิเวศนำหน้าในกลุ่มผู้ผลิตรถยนต์ดั้งเดิม
\end{itemize}

\booksection{ใช้โทรศัพท์แทนหน้าจอรถยนต์ — ใช้ได้ไหม?}

ถ้าหน้าจอรถยนต์ของคุณเก่าเกินไปหรือใช้ยาก สามารถใช้โทรศัพท์แทนได้เลยไหม? คำตอบคือ ได้ แต่ต้องดูสถานการณ์

\begin{itemize}
  \item \textbf{Motorola Ready For / โหมดเดสก์ท็อป}: โทรศัพท์ Motorola, Samsung และ Huawei บางรุ่นรองรับ "โหมดเดสก์ท็อป" เสียบกับจอภาพก็กลายเป็นคอมพิวเตอร์เครื่องง่ายๆ ตามทฤษฎีแล้วเสียบกับหน้าจอรถยนต์ก็ใช้ได้ แต่ความเข้ากันได้จริงนั้นแตกต่างกันไป และขาดการโต้ตอบเฉพาะสำหรับรถยนต์ (เช่น ปุ่มบนพวงมาลัย)
  \item \textbf{ขั้วจับโทรศัพท์ vs หน้าจอรถยนต์เดิม}: หากหน้าจอรถยนต์เดิมค้าง แผนที่ไม่อัปเดต การใช้ขั้วจับโทรศัพท์ที่แข็งแรง + นำทางด้วยโทรศัพท์ถือเป็นทางเลือกที่ใช้งานได้จริง แต่ต้องระวัง: (1)หน้าจอโทรศัพท์เล็ก ขับรถแล้วมองไม่ชัด (2)แดดจัดในหน้าร้อนจะทำให้โทรศัพท์ร้อนจัด ลดความเร็วหรือดับได้ (3)นำทางนานๆ ใช้พลังงานเร็ว ต้องใช้ที่ชาร์จในรถ
  \item \textbf{คำแนะนำ}: ระยะสั้นหรือใช้รถชั่วคราว ขั้วจับโทรศัพท์ก็เพียงพอ การเดินทางสัญจรประจำวันหรือระยะไกล แนะนำให้อัปเกรดเป็นหน้าจอที่รองรับ CarPlay/HiCar หรือติดตั้งหน้าจอรถยนต์แบบ Aftermarket ทั้งประสบการณ์และความปลอดภัยจะดีกว่า
\end{itemize}

\begin{tcbwarn}{คำเตือนด้านความปลอดภัย}
ไม่ว่าจะใช้วิธีใด การใช้โทรศัพท์ขณะขับรถถือเป็นการฝ่าฝืนกฎจราจร ขั้วจับโทรศัพท์ใช้สำหรับ “ดูแผนที่” ได้เท่านั้น ไม่สามารถใช้เพื่อ “กดใช้งานโทรศัพท์” การตั้งค่าจุดหมายปลายทาง สลับเพลง ฯลฯ โปรดอย่าทำขณะขับรถ หรือใช้คำสั่งเสียง

\end{tcbwarn}

\booksection{เคล็ดลับที่ใช้งานได้จริง}

\begin{itemize}
  \item \textbf{การเชื่อมต่อ CarPlay แบบไร้สายไม่เสถียร}: (1)ปิด Bluetooth ของหน้าจอรถยนต์แล้วจับคู่ใหม่ (2)ลบประวัติอุปกรณ์ที่จับคู่ไว้ในหน้าจอรถยนต์แล้วเชื่อมต่อใหม่ (3)ตรวจสอบว่าโทรศัพท์เชื่อมต่อกับอุปกรณ์ Bluetooth หลายตัวพร้อมกันหรือไม่ซึ่งอาจทำให้เกิดความขัดแย้ง (4)อัปเกรดเฟิร์มแวร์ของหน้าจอรถยนต์และระบบปฏิบัติการของโทรศัพท์เป็นเวอร์ชันล่าสุด (5)หากยังไม่ได้ผล ให้เปลี่ยนไปใช้ CarPlay แบบมีสาย จะเสถียรกว่า
  \item \textbf{ปริมาณข้อมูลอินเทอร์เน็ตของหน้าจอรถยนต์หมด}: (1)เปิดฮอตสปอตจากโทรศัพท์ แล้วให้หน้าจอรถยนต์เชื่อมต่อ Wi-Fi จากฮอตสปอตของโทรศัพท์ (2)ในการตั้งค่าหน้าจอรถยนต์ ให้สลับ App ที่ "ใช้ข้อมูลมาก" (เช่น วิดีโอ เพลงออนไลน์) ไปใช้ Wi-Fi (3)ล้างแคชแผนที่ออฟไลน์ของหน้าจอรถยนต์เป็นประจำ
  \item \textbf{การเปิดโหมดขับรถอัตโนมัติ}: (1)ตั้งค่า "โหมดโฟกัส" ของ iOS ให้ "เปิดอัตโนมัติเมื่อเชื่อมต่อกับ Bluetooth ของรถยนต์" เพื่อบล็อกการแจ้งเตือนที่ไม่เร่งด่วน (2)"โหมดขับรถ" ของ Android หรือ App ของบุคคลที่สาม (เช่น Tasker) สามารถตั้งค่าให้ปิดเสียงอัตโนมัติเมื่อเชื่อมต่อ Bluetooth และตอบกลับข้อความอัตโนมัติ (3)โทรศัพท์ Huawei สามารถตั้งค่า "การชำระเงินอัจฉริยะ" เพื่อสลับไปเล่นเพลงบนหน้าจอรถยนต์โดยอัตโนมัติเมื่อขึ้นรถ
\end{itemize}

\begin{tcbinfo}{คำแนะนำเล็กๆ น้อยๆ เวลาเลือกรถ}
หากคุณสนใจประสบการณ์การใช้งานหน้าจอรถยนต์ เวลาเลือกรถต้องไปทดสอบที่โชว์รูมจริงๆ: (1)นำโทรศัพท์ของคุณไปทดลองเชื่อมต่อ CarPlay/HiCar ว่าเสถียรหรือไม่ (2)ทดสอบความเร็วในการตอบสนองและอัตราการจดจำคำสั่งของผู้ช่วยเสียง (3)ดูว่า App Store มี App ที่คุณใช้บ่อยหรือไม่ (4)ติดตามความถี่ในการอัปเดต OTA ในอนาคต — หน้าจอรถยนต์ที่ดีเกิดจากการพัฒนาซอฟต์แวร์อย่างต่อเนื่อง

\end{tcbinfo}

\begin{tcbpractice}{ลงมือทำ: ปรับปรุงประสบการณ์การเชื่อมต่อหน้าจอรถยนต์ของคุณ}
\begin{enumerate}
  \item ยืนยันวิธีการเชื่อมต่อที่โทรศัพท์และหน้าจอรถยนต์ของคุณรองรับ (CarPlay / HiCar / Android Auto / CarLife)
  \item ล้างรายการจับคู่ Bluetooth ของหน้าจอรถยนต์ โดยเก็บเฉพาะโทรศัพท์ที่ใช้บ่อยไว้ เพื่อลดความขัดแย้งในการเชื่อมต่อ
  \item ตั้งค่า "โหมดโฟกัสขณะขับรถ" ให้เปิดอัตโนมัติเมื่อเชื่อมต่อกับ Bluetooth ของรถยนต์
  \item ดาวน์โหลด App ของแบรนด์รถยนต์ (เช่น ฮงเหมิงจื้อซิง, Li Auto, NIO App) เพื่อดูว่าสามารถควบคุมแอร์และตรวจสอบสภาพรถจากระยะไกลได้หรือไม่
\end{enumerate}

\end{tcbpractice}

\begin{tcbsummary}{สรุปบทนี้}
\begin{itemize}
  \item สามเส้นทางหลักของการเชื่อมต่อหน้าจอรถยนต์: การมิเรอร์หน้าจอ (CarPlay ฯลฯ), การผสานรวมอย่างลึกซึ้ง (ห้องโดยสาร HarmonyOS ฯลฯ), และการผสานรวมอย่างลึกของ CarPlay 2.0
  \item ระบบนิเวศของ CarPlay กว้างขวางที่สุด; HiCar ให้ประสบการณ์ที่ดีในประเทศ; Android Auto ต้องปรับแต่งในประเทศ; CarLife เหมาะสำหรับการติดตั้งเพิ่มบนรถเก่า
  \item ซูเปอร์เดสก์ท็อป การส่งต่อแอปพลิเคชัน และการใช้บัญชีร่วมกันของห้องโดยสาร HarmonyOS เป็นตัวแทนของทิศทางการผสานรวมอย่างลึกซึ้ง
  \item NOMI ของ NIO, หลายหน้าจอของ Li Auto, ระบบนิเวศบุคคล-รถ-บ้านของ Xiaomi, และหน้าจอหมุนได้ของ BYD แต่ละแบรนด์มีจุดเด่นเฉพาะตัว
  \item การใช้โทรศัพท์แทนหน้าจอรถยนต์เหมาะสำหรับสถานการณ์ชั่วคราวเท่านั้น ในระยะยาวแนะนำให้อัปเกรดหรือติดตั้งหน้าจอรถยนต์เพิ่มเติม
  \item เคล็ดลับที่ใช้งานได้จริง: การเชื่อมต่อ CarPlay ใหม่ ใช้ฮอตสปอตโทรศัพท์เสริมปริมาณข้อมูล และใช้ Bluetooth เปิดโหมดขับรถ
\end{itemize}

\end{tcbsummary}
